\documentclass[aspectratio=169, 10pt]{beamer}

% ==============================================================================
% CẤU HÌNH GÓI & MÔI TRƯỜNG BIÊN DỊCH TIẾNG VIỆT ĐA NỀN TẢNG
% ==============================================================================
\usepackage{iftex}
\ifPDFTeX
  \usepackage[utf8]{inputenc}
  \usepackage[vietnamese]{babel}
  \usepackage[T5]{fontenc}
  \usepackage{lmodern}
\else
  \usepackage{fontspec}
  \usepackage[vietnamese]{babel}
\fi

\usepackage{tikz}
\usetikzlibrary{shapes.geometric, arrows.meta, positioning, calc, backgrounds, fit, patterns}
\usepackage{xcolor}
\usepackage{listings}
\usepackage{booktabs}
\usepackage{array}
\usepackage{colortbl}
\usepackage{tcolorbox}
\tcbuselibrary{skins, listings, breakable}

% ==============================================================================
% BẢNG MÀU CHUẨN THIẾT KẾ ĐỒNG BỘ THEO MẪU BÁO CÁO (NAVY & ORANGE ACCENT)
% ==============================================================================
\definecolor{NavyHeader}{RGB}{26, 43, 70}      % #1A2B46: Xanh Navy chuẩn thanh tiêu đề
\definecolor{AccentOrange}{RGB}{229, 126, 33}  % #E57E21: Cam vát góc chuẩn thương hiệu
\definecolor{NavyTitle}{RGB}{17, 34, 51}       % #112233: Tiêu đề slide chính
\definecolor{TextDark}{RGB}{44, 62, 80}        % #2C3E50: Nội dung văn bản
\definecolor{CodeBg}{RGB}{248, 249, 250}       % #F8F9FA: Nền hộp code nhẹ
\definecolor{CodeBorder}{RGB}{226, 232, 240}   % #E2E8F0: Viền xám hộp code
\definecolor{SuccessGreen}{RGB}{46, 125, 50}   % #2E7D32: Xanh lá cây thành công
\definecolor{ErrorRed}{RGB}{198, 40, 40}       % #C62828: Đỏ cảnh báo lỗi
\definecolor{CardBg}{RGB}{245, 247, 250}       % Nền khối nổi bật
\definecolor{TableHeadBg}{RGB}{237, 242, 247}  % Nền đầu bảng
\definecolor{TableAltBg}{RGB}{248, 250, 252}   % Dòng so le bảng

% ==============================================================================
% ĐỊNH DẠNG BEAMER: THANH TIÊU ĐỀ TRÊN CÙNG CHUẨN HÌNH 1 (VIỀN VÁT GÓC ĐỘC LẬP)
% ==============================================================================
\setbeamertemplate{navigation symbols}{}
\setbeamersize{text margin left=0.85cm, text margin right=0.85cm}
\setbeamercolor{background canvas}{bg=white}
\setbeamercolor{normal text}{fg=TextDark}

% Header chuẩn xác 100% theo Ảnh 1:
% Nền trắng bên trái -> Dải vạch cam chéo -> Khoảng hở trắng chéo -> Thanh Navy vát chéo
\setbeamertemplate{headline}{%
  \begin{tikzpicture}[remember picture, overlay]
    % 1. Dải vát màu Cam (Parallelogram độc lập trên nền trắng)
    \fill[AccentOrange] 
      ([xshift=2.85cm]current page.north west) -- 
      ([xshift=3.13cm]current page.north west) -- 
      ([xshift=4.08cm, yshift=-0.85cm]current page.north west) -- 
      ([xshift=3.80cm, yshift=-0.85cm]current page.north west) -- cycle;
      
    % 2. Thanh Navy chính vát chéo song song (cách dải cam 1 khoảng hở trắng chuẩn)
    \fill[NavyHeader] 
      ([xshift=3.35cm]current page.north west) -- 
      (current page.north east) -- 
      ([yshift=-0.85cm]current page.north east) -- 
      ([xshift=4.30cm, yshift=-0.85cm]current page.north west) -- cycle;
      
    % 3. Chữ tiêu đề trên nền thanh Navy
    \node[anchor=center, text=white, font=\small\bfseries] at ([xshift=0.3cm, yshift=-0.425cm]current page.north) {Kiến trúc ứng dụng Flutter -- Bố cục \& Không gian};
    % 4. Số trang ở góc phải
    \node[anchor=east, text=white, font=\small\bfseries] at ([xshift=-1.0cm, yshift=-0.425cm]current page.north east) {\insertframenumber{} / \inserttotalframenumber};
  \end{tikzpicture}%
  \vspace{0.85cm}%
}

\setbeamertemplate{footline}{}

% ==============================================================================
% CẤU HÌNH ĐỊNH DẠNG CODE DART (LISTINGS)
% ==============================================================================
\lstdefinelanguage{Dart}{
  morekeywords={
    abstract, as, assert, async, await, break, case, catch, class, const, continue,
    default, deferred, do, dynamic, else, enum, export, extends, external, factory,
    false, final, finally, for, get, if, implements, import, in, is, library, new,
    null, operator, part, rethrow, return, set, static, super, switch, sync, this,
    throw, true, try, typedef, var, void, while, with, yield,
    Widget, BuildContext, State, StatefulWidget, StatelessWidget, Key,
    Row, Column, Stack, Positioned, Expanded, Flexible, SingleChildScrollView,
    ListView, Wrap, FittedBox, LayoutBuilder, ConstrainedBox, IntrinsicHeight, Spacer, Text, Container,
    BoxConstraints, MainAxisAlignment, CrossAxisAlignment, FlexFit, TextOverflow,
    Color, Colors, EdgeInsets, TextStyle, FontWeight, BorderRadius, BoxDecoration,
    Scaffold, AppBar, Center, SizedBox, Icon, Icons, CircleAvatar, SnackBar,
    Axis, Alignment, Border, BoxShadow, Offset, MediaQuery
  },
  sensitive=true,
  morecomment=[l]{//},
  morecomment=[s]{/*}{*/},
  morestring=[b]',
  morestring=[b]"
}

\lstset{
  language=Dart,
  backgroundcolor=\color{CodeBg},
  basicstyle=\ttfamily\scriptsize\color{TextDark},
  keywordstyle=\bfseries\color{NavyTitle},
  commentstyle=\itshape\color{gray!80!black},
  stringstyle=\color{AccentOrange!90!black},
  numbers=left,
  numberstyle=\tiny\color{gray!60!black},
  numbersep=6pt,
  frame=single,
  rulecolor=\color{CodeBorder},
  breaklines=true,
  showstringspaces=false,
  tabsize=2,
  xleftmargin=12pt,
  framexleftmargin=10pt
}

% Lệnh tiện ích tiêu đề slide chuẩn
\newcommand{\SlideHeader}[2]{%
  \noindent{\textcolor{AccentOrange}{\textbf{\footnotesize #1}}}\\[0.08cm]
  \noindent{\Large \textbf{\textcolor{NavyTitle}{#2}}}\\[0.22cm]
}

\begin{document}

% ==============================================================================
% TRANG 1: TRANG BÌA BÁO CÁO CHUYÊN ĐỀ (ĐÃ BỎ CÁC CHI TIẾT THỪA THEO YÊU CẦU ẢNH 2)
% ==============================================================================
\begin{frame}
  \vspace{0.75cm}
  \centering
  {\textcolor{AccentOrange}{\textbf{\footnotesize BÁO CÁO CHUYÊN ĐỀ KỸ THUẬT}}}\\[0.32cm]
  {\Huge \textbf{\textcolor{NavyTitle}{BỐ CỤC FLUTTER: FLEX \& SPACE}}}\\[0.20cm]
  {\LARGE \textbf{\textcolor{NavyTitle}{Tổ Chức Bố Cục Linh Hoạt \& Xử Lý Không Gian}}}\\[0.40cm]
  {\color{AccentOrange}\rule{5.2cm}{1.3pt}}\\[0.50cm]

  {\large \textbf{Chủ đề 02:} Main Axis, Cross Axis -- Expanded, Flexible, Stack \& Positioned}\\[0.18cm]
  {\normalsize Toàn diện 5 phương pháp xử lý Overflow \& Giải pháp khắc phục triệt để trong dự án}\\[0.45cm]

  {\small Đầy đủ cấu trúc ứng dụng 7 thành phần \& Demo phản ứng khi kích thước / nội dung thay đổi}
\end{frame}

% ==============================================================================
% TRANG 2: TỔNG QUAN NỘI DUNG BÁO CÁO (AGENDA)
% ==============================================================================
\begin{frame}
  \SlideHeader{TỔNG QUAN CHUYÊN ĐỀ}{Nội Dung Báo Cáo (15 Phút)}

  \begin{columns}[T]
    % Cột I
    \begin{column}{0.48\textwidth}
      {\textcolor{AccentOrange}{\textbf{PHẦN I: LÝ THUYẾT \& 5 CÁCH XỬ LÝ OVERFLOW}}}\\[0.18cm]

      \textbf{1. Đặt vấn đề \& Thách thức không gian}\\
      {\footnotesize Vấn nạn vỡ màn hình, code cứng kích thước \& vạch sọc vàng-đen.}\\[0.14cm]

      \textbf{2. Cốt lõi: Nguyên lý Box Constraints}\\
      {\footnotesize Khẩu quyết: Constraints down, Sizes up, Parent sets position.}\\[0.14cm]

      \textbf{3. Trục toạ độ: Main Axis \& Cross Axis}\\
      {\footnotesize Phân định hướng trục và các chế độ căn chỉnh trên Row \& Column.}\\[0.14cm]

      \textbf{4. Phân chia không gian: Expanded vs Flexible}\\
      {\footnotesize So sánh Tight Fit (ép lấp đầy) vs Loose Fit (co giãn tự nhiên).}\\[0.14cm]

      \textbf{5. Bố cục xếp chồng: Stack \& Positioned}\\
      {\footnotesize Quản lý trục Z và định vị toạ độ neo chính xác.}\\[0.14cm]

      \textbf{6. Nguyên nhân gây lỗi Overflow (RenderFlex)}\\
      {\footnotesize Cơ chế vượt kích thước tối đa và 3 tình huống kích hoạt lỗi.}\\[0.14cm]

      \textbf{7. Chi tiết 5 cách xử lý Overflow chuẩn giáo trình}\\
      {\footnotesize Cuộn trang, Phân chia linh hoạt, Wrap, Ellipsis, FittedBox.}
    \end{column}

    % Cột II
    \begin{column}{0.48\textwidth}
      {\textcolor{AccentOrange}{\textbf{PHẦN II: THỰC THI DEMO \& CODE DỰ ÁN}}}\\[0.18cm]

      \textbf{8. Cấu trúc ứng dụng thực tế 7 thành phần}\\
      {\footnotesize Thiết kế màn hình đăng nhập ngân hàng chuẩn Clean UI.}\\[0.14cm]

      \textbf{9. Minh họa mã nguồn: Căn chỉnh Row \& Column}\\
      {\footnotesize Đoạn code áp dụng spaceBetween và căn lề trục phụ.}\\[0.14cm]

      \textbf{10. Minh họa mã nguồn: Expanded \& Nút cố định}\\
      {\footnotesize Phân bổ tỷ lệ hàng nút Login co dãn và Face ID 52x52.}\\[0.14cm]

      \textbf{11. Nhấn mạnh giải pháp của tôi: Tràn ngang (Ellipsis)}\\
      {\footnotesize Code AccountCard trị tên siêu dài bằng Expanded + Ellipsis.}\\[0.14cm]

      \textbf{12. Nhấn mạnh giải pháp của tôi: Bộ 4 trị tràn dọc}\\
      {\footnotesize Code LayoutBuilder + ScrollView + ConstrainedBox chống tràn bàn phím.}\\[0.14cm]

      \textbf{13. Nhấn mạnh giải pháp của tôi: Flexible \& Responsive}\\
      {\footnotesize Code nút tùy chọn phụ fit loose và khống chế maxWidth 440px.}\\[0.14cm]

      \textbf{14. So sánh kỹ thuật, Kiểm chứng Demo \& Tổng kết}\\
      {\footnotesize Bảng đối sánh 5 giải pháp, kiểm thử trực quan \& Q\&A.}
    \end{column}
  \end{columns}
\end{frame}

% ==============================================================================
% TRANG 3: ĐẶT VẤN ĐỀ - HIỆN TRẠNG & THÁCH THỨC
% ==============================================================================
\begin{frame}
  \SlideHeader{ĐẶT VẤN ĐỀ}{Hiện Trạng \& Thách Thức Khi Xử Lý Bố Cục Không Gian}

  \textbf{Vấn nạn vỡ giao diện trên đa dạng màn hình (Screen Fragmentation)}\\
  {\small Thói quen viết code gán cứng kích thước tuyệt đối (\texttt{width: 380, height: 220}) khiến giao diện chỉ chạy đẹp trên thiết bị thử nghiệm, nhưng lập tức bị biến dạng, méo mó trên các dòng máy khác nhau.}\\[0.22cm]

  \textbf{Ác mộng vạch sọc vàng-đen: "A RenderFlex overflowed by xxx pixels"}\\
  {\small Xuất hiện ngay khi người dùng nhập chuỗi văn bản dài hơn dự kiến, hoặc khi bàn phím ảo xuất hiện chiếm từ 250px -- 350px không gian dọc, khiến luồng người dùng bị gián đoạn.}\\[0.22cm]

  \textbf{Lạm dụng lồng ghép Row, Column không kiểm soát (Spaghetti Tree)}\\
  {\small Không hiểu rõ mối quan hệ ràng buộc kích thước (Constraints) giữa Widget cha và Widget con, dẫn đến các lỗi kinh điển như \textit{Unbounded height/width} gây crash ứng dụng trong pha Render.}\\[0.22cm]

  \textbf{Ứng dụng hoàn toàn thiếu khả năng thích ứng linh hoạt (Non-responsive)}\\
  {\small Khi người dùng xoay ngang màn hình (Landscape) hoặc chuyển sang chế độ đa nhiệm chia đôi màn hình (Split-screen), bố cục bị cắt cụt và không thể tự phân bổ lại không gian.}
\end{frame}

% ==============================================================================
% TRANG 4: CƠ SỞ LÝ THUYẾT - NGUYÊN LÝ BOX CONSTRAINTS
% ==============================================================================
\begin{frame}
  \SlideHeader{CƠ SỞ LÝ THUYẾT}{Nguyên Lý Ràng Buộc Hộp: "Box Constraints"}

  \begin{columns}[T]
    \begin{column}{0.54\textwidth}
      \textbf{Khẩu quyết 3 bước kinh điển trong Flutter Layout:}
      \begin{enumerate}
        \item \textbf{Constraints go down (Ràng buộc đi xuống):}\\
        Widget cha truyền các giới hạn về độ rộng và chiều cao cực tiểu/cực đại (\texttt{minWidth, maxWidth, minHeight, maxHeight}) xuống widget con.
        \item \textbf{Sizes go up (Kích thước báo cáo lên):}\\
        Widget con tự chọn kích thước của chính nó dựa trên ràng buộc được cấp, rồi gửi kích thước đó ngược lên cha.
        \item \textbf{Parent sets position (Cha định vị toạ độ):}\\
        Widget cha nhận kích thước từ các con và quyết định toạ độ $(x, y)$ đặt con lên màn hình.
      \end{enumerate}
      \vspace{0.1cm}
      \textbf{3 Phân loại ràng buộc cốt lõi:}
      \begin{itemize}
        \item \textbf{Tight}: \texttt{min == max} (Ép buộc con nhận chính xác kích thước).
        \item \textbf{Loose}: \texttt{min = 0, max > 0} (Con tự do từ 0 đến max).
        \item \textbf{Unbounded}: \texttt{max = infinity} (Không giới hạn kích thước).
      \end{itemize}
    \end{column}

    \begin{column}{0.44\textwidth}
      \centering
      \begin{tikzpicture}[node distance=1.1cm, font=\scriptsize]
        % Node Cha
        \node[draw=NavyTitle, fill=NavyHeader!10, rounded corners=6pt, text width=4.0cm, align=center, line width=1.1pt] (parent) {
          \textbf{\textcolor{NavyTitle}{WIDGET CHA (Parent)}}\\
          {\tiny Nắm giữ không gian màn hình}
        };

        % Node Con
        \node[draw=AccentOrange, fill=AccentOrange!10, rounded corners=6pt, text width=4.0cm, align=center, below=of parent, line width=1.1pt] (child) {
          \textbf{\textcolor{AccentOrange}{WIDGET CON (Child)}}\\
          {\tiny Quyết định size trong giới hạn}
        };

        % Mũi tên 1: Ràng buộc đi xuống
        \draw[->, >=stealth, line width=1.2pt, color=NavyTitle] ([xshift=-1.0cm]parent.south) -- ([xshift=-1.0cm]child.north)
          node[midway, left, text width=2.2cm, align=right] {\tiny \textbf{1. Constraints}\\min/max $(w, h)$};

        % Mũi tên 2: Kích thước đi lên
        \draw[->, >=stealth, line width=1.2pt, color=AccentOrange] ([xshift=1.0cm]child.north) -- ([xshift=1.0cm]parent.south)
          node[midway, right, text width=2.0cm, align=left] {\tiny \textbf{2. Size}\\$(w, h)$ cụ thể};

        % Hộp Toạ độ 3
        \node[draw=SuccessGreen, fill=SuccessGreen!10, rounded corners=4pt, text width=3.8cm, align=center, below=0.6cm of child] (pos) {
          \textbf{\textcolor{SuccessGreen}{3. Parent sets Position $(x, y)$}}\\
          {\tiny Định vị con trên canvas đồ họa}
        };
        \draw[->, >=stealth, dashed, line width=1.0pt, color=SuccessGreen] (child.south) -- (pos.north);
      \end{tikzpicture}
    \end{column}
  \end{columns}
\end{frame}

% ==============================================================================
% TRANG 5: CƠ SỞ LÝ THUYẾT - MAIN AXIS & CROSS AXIS
% ==============================================================================
\begin{frame}
  \SlideHeader{CƠ SỞ LÝ THUYẾT}{Trục Chính (Main Axis) \& Trục Phụ (Cross Axis)}

  \begin{columns}[T]
    \begin{column}{0.48\textwidth}
      \textbf{Row (Bố cục theo hàng ngang):}
      \begin{itemize}
        \item \textbf{Main Axis (Trục chính):} Nằm ngang ($X$).\\
        Thuộc tính \texttt{mainAxisAlignment} quyết định cách chia sẻ khoảng trống từ mép trái sang mép phải.
        \item \textbf{Cross Axis (Trục phụ):} Thẳng đứng ($Y$).\\
        Thuộc tính \texttt{crossAxisAlignment} căn chỉnh độ cao các thành phần (\texttt{center, start, end, stretch}).
      \end{itemize}

      \vspace{0.25cm}
      \textbf{Column (Bố cục theo cột dọc):}
      \begin{itemize}
        \item \textbf{Main Axis (Trục chính):} Thẳng đứng ($Y$).\\
        Điều phối khoảng cách từ trên đỉnh xuống chân cột.
        \item \textbf{Cross Axis (Trục phụ):} Nằm ngang ($X$).\\
        Căn mép trái, giữa hoặc kéo dãn (\texttt{stretch}) theo bề rộng cột.
      \end{itemize}
    \end{column}

    \begin{column}{0.48\textwidth}
      \textbf{Các giá trị MainAxisAlignment chuẩn mực:}
      \begin{itemize}
        \item \textbf{start / center / end:}\\
        Dồn toàn bộ phần tử về đầu, chính giữa hoặc cuối trục.
        \item \textbf{spaceBetween:}\\
        Đẩy phần tử đầu và cuối sát mép ngoài, chia đều khoảng trống ở giữa.
        \item \textbf{spaceAround:}\\
        Chia đều không gian xung quanh mỗi phần tử (khoảng đệm ở biên bằng $1/2$ ở giữa).
        \item \textbf{spaceEvenly:}\\
        Tất cả khoảng cách giữa các phần tử và mép ngoài đều bằng nhau chằn chặn.
      \end{itemize}

      \vspace{0.15cm}
      \begin{tcolorbox}[colback=CardBg, colframe=NavyHeader!30, arc=4pt, boxrule=0.8pt, left=6pt, right=6pt, top=4pt, bottom=4pt]
        {\scriptsize \textbf{Lưu ý vàng:} Row và Column tự cấp ràng buộc \textbf{Unbounded} trên trục chính của chúng. Do đó không thể đặt một widget co giãn tự do không giới hạn vào bên trong mà không có bộ điều phối.}
      \end{tcolorbox}
    \end{column}
  \end{columns}
\end{frame}

% ==============================================================================
% TRANG 6: CƠ SỞ LÝ THUYẾT - EXPANDED VS FLEXIBLE
% ==============================================================================
\begin{frame}
  \SlideHeader{CƠ SỞ LÝ THUYẾT}{Phân Tách Không Gian: Expanded vs Flexible}

  \begin{columns}[T]
    \begin{column}{0.48\textwidth}
      \textbf{1. Expanded (Ép lấp đầy -- Tight Fit):}
      \begin{itemize}
        \item Mặc định sử dụng \texttt{fit: FlexFit.tight}.
        \item Ép buộc widget con \textbf{phải lấp đầy 100\%} phần không gian được phân bổ, bất kể kích thước nội tại của con lớn hay bé.
        \item Bản chất \texttt{Expanded} chính là một cách viết ngắn gọn của \texttt{Flexible(fit: FlexFit.tight)}.
        \item Ứng dụng: Căng nút bấm hết chiều ngang, mở rộng ô tìm kiếm, bọc chuỗi văn bản dài.
      \end{itemize}

      \vspace{0.2cm}
      \textbf{2. Flexible (Co giãn tự do -- Loose Fit):}
      \begin{itemize}
        \item Mặc định sử dụng \texttt{fit: FlexFit.loose}.
        \item Cho phép con giữ nguyên kích thước tự nhiên, nhưng đặt \textbf{giới hạn trần} không được vượt quá phần không gian được chia.
        \item Ứng dụng: Giữ chữ nhỏ tự nhiên, không ép dãn vô nghĩa.
      \end{itemize}
    \end{column}

    \begin{column}{0.48\textwidth}
      \textbf{3. Hệ số \texttt{flex} (Trọng số phân bổ):}
      \begin{itemize}
        \item Không gian còn trống trên trục chính sẽ được chia thành tổng các phần bằng $\sum \text{flex}$.
        \item Mỗi widget nhận về phần không gian tỷ lệ:
        $$\text{Không gian con} = \text{Không gian trống} \times \frac{\text{flex}_i}{\sum \text{flex}}$$
      \end{itemize}

      \vspace{0.2cm}
      \textbf{Sơ đồ phân bổ trực quan (Row với flex 2 : 1):}
      \begin{tikzpicture}[font=\scriptsize]
        \draw[fill=NavyHeader!15, draw=NavyHeader, line width=1pt] (0, 0) rectangle (3.2, 0.7);
        \node at (1.6, 0.35) {\textbf{\textcolor{NavyTitle}{Expanded (flex: 2)}}};

        \draw[fill=AccentOrange!20, draw=AccentOrange, line width=1pt] (3.3, 0) rectangle (4.9, 0.7);
        \node at (4.1, 0.35) {\textbf{\textcolor{AccentOrange}{Flex: 1}}};

        \node[above] at (1.6, 0.7) {\tiny Chiếm 2/3 không gian trống};
        \node[above] at (4.1, 0.7) {\tiny Chiếm 1/3};
      \end{tikzpicture}
    \end{column}
  \end{columns}
\end{frame}

% ==============================================================================
% TRANG 7: CƠ SỞ LÝ THUYẾT - STACK & POSITIONED
% ==============================================================================
\begin{frame}
  \SlideHeader{CƠ SỞ LÝ THUYẾT}{Bố Cục Xếp Lớp Đa Chiều: Stack \& Positioned}

  \begin{columns}[T]
    \begin{column}{0.52\textwidth}
      \textbf{Quản lý không gian theo trục $Z$ (Trục chiều sâu):}
      \begin{itemize}
        \item Không bố trí nối tiếp tuần tự như Row/Column mà xếp chồng các lớp giao diện lên nhau.
        \item \textbf{Quy tắc thứ tự vẽ (Painter Order):} Widget khai báo trước nằm ở lớp đáy (Background); widget khai báo sau cùng nằm ở trên cùng (Foreground).
      \end{itemize}

      \vspace{0.15cm}
      \textbf{Positioned -- Định vị toạ độ tuyệt đối:}
      \begin{itemize}
        \item Là widget con đặc thù chỉ dùng bên trong \texttt{Stack}.
        \item Cố định chính xác vị trí qua các thuộc tính: \texttt{top}, \texttt{bottom}, \texttt{left}, \texttt{right}, kèm theo \texttt{width}, \texttt{height}.
      \end{itemize}

      \vspace{0.15cm}
      \textbf{Non-positioned Children (Thành phần phi toạ độ):}
      \begin{itemize}
        \item Được căn chỉnh tự động dựa trên thuộc tính \texttt{Stack.alignment} (Mặc định: \texttt{AlignmentDirectional.topStart}).
      \end{itemize}
    \end{column}

    \begin{column}{0.44\textwidth}
      \centering
      \textbf{Minh hoạ cấu trúc xếp lớp trục $Z$:}\\[0.25cm]
      \begin{tikzpicture}[scale=0.8, font=\scriptsize]
        % Lớp 0: Nền Card
        \draw[fill=gray!20, draw=gray!60, line width=1pt] (0, 0) -- (3.5, 0) -- (4.2, 1.2) -- (0.7, 1.2) -- cycle;
        \node at (2.1, 0.6) {\textcolor{gray!80!black}{\textbf{Lớp 0: Thẻ chứa (Container)}}};

        % Lớp 1: Avatar
        \draw[fill=NavyHeader!30, draw=NavyHeader, line width=1.1pt] (0.4, 0.9) -- (3.9, 0.9) -- (4.6, 2.1) -- (1.1, 2.1) -- cycle;
        \node at (2.5, 1.5) {\textbf{\textcolor{NavyTitle}{Lớp 1: CircleAvatar}}};

        % Lớp 2: Positioned Badge
        \draw[fill=SuccessGreen!40, draw=SuccessGreen, line width=1.2pt] (0.8, 1.8) -- (4.3, 1.8) -- (5.0, 3.0) -- (1.5, 3.0) -- cycle;
        \node at (2.9, 2.4) {\textbf{\textcolor{SuccessGreen}{Lớp 2: Positioned (Online Badge)}}};

        % Trục Z
        \draw[->, >=stealth, line width=1.3pt, color=AccentOrange] (-0.3, -0.2) -- (0.5, 3.3) node[above] {\textbf{\textcolor{AccentOrange}{Trục Z}}};
      \end{tikzpicture}
    \end{column}
  \end{columns}
\end{frame}

% ==============================================================================
% TRANG 8: CƠ SỞ LÝ THUYẾT - NGUYÊN NHÂN LỖI OVERFLOW
% ==============================================================================
\begin{frame}
  \SlideHeader{CƠ SỞ LÝ THUYẾT}{Nguyên Nhân \& Cơ Chế Gây Ra Lỗi Overflow}

  \begin{columns}[T]
    \begin{column}{0.54\textwidth}
      \textbf{Bản chất kỹ thuật của RenderFlex Overflow:}
      \begin{itemize}
        \item Xảy ra trong pha bố cục (\textit{Layout Phase}) của Flutter Rendering Engine.
        \item Khi tổng kích thước các widget con trên trục chính lớn hơn kích thước tối đa (\texttt{maxConstraints}) mà cha cho phép:
        $$\Delta_{\text{overflow}} = \sum \text{Size}_{\text{children}} - \text{MaxConstraint}_{\text{parent}} > 0$$
      \end{itemize}

      \vspace{0.15cm}
      \textbf{3 Tình huống kích hoạt lỗi thường gặp nhất:}
      \begin{enumerate}
        \item \textbf{Tràn ngang do văn bản:}\\
        Đặt \texttt{Text} trong \texttt{Row} mà không bọc \texttt{Expanded}. Chuỗi dài đẩy vượt mép màn hình.
        \item \textbf{Tràn dọc do bàn phím ảo:}\\
        Bàn phím chiếm diện tích làm giảm chiều cao màn hình, trong khi \texttt{Column} không có cơ chế cuộn.
        \item \textbf{Lỗi lồng ghép Unbounded:}\\
        Đặt \texttt{ListView} bên trong \texttt{Column} mà không chỉ định chiều cao hoặc bọc \texttt{Expanded}.
      \end{enumerate}
    \end{column}

    \begin{column}{0.44\textwidth}
      \centering
      \begin{tcolorbox}[colback=red!5, colframe=ErrorRed, arc=6pt, boxrule=1.2pt, left=8pt, right=8pt, top=8pt, bottom=8pt]
        \centering
        \textbf{\textcolor{ErrorRed}{\small CẢNH BÁO DEBUG CONSOLE:}}\\[0.15cm]
        {\scriptsize \ttfamily A RenderFlex overflowed by 142.0 pixels on the right.}\\[0.25cm]
        \begin{tikzpicture}
          % Thanh sọc vằn vàng-đen đặc trưng Flutter
          \draw[line width=12pt, pattern=north east lines, pattern color=black, draw=yellow!90!black] (0, 0) -- (4.2, 0);
        \end{tikzpicture}\\[0.15cm]
        {\tiny Vạch sọc vàng-đen xuất hiện báo hiệu giao diện đã bị xé rách không gian!}
      \end{tcolorbox}
    \end{column}
  \end{columns}
\end{frame}

% ==============================================================================
% TRANG 9: CƠ SỞ LÝ THUYẾT - 5 CÁCH XỬ LÝ OVERFLOW (PHẦN 1)
% ==============================================================================
\begin{frame}
  \SlideHeader{CƠ SỞ LÝ THUYẾT}{Chi Tiết Các Cách Xử Lý Overflow (Phần 1)}

  \textbf{1. Sử dụng cuộn trang (\texttt{SingleChildScrollView} hoặc \texttt{ListView}):}
  \begin{itemize}
    \item \textbf{Nguyên tắc:} Khi nội dung bên trong vượt quá chiều cao hoặc chiều rộng của màn hình, thay vì cố gắng ép hiển thị trong một khung cố định, ta bọc nó vào một widget hỗ trợ cuộn.
    \item \textbf{Áp dụng:} Thường dùng cho các form nhập liệu dài, trang cài đặt, hoặc chi tiết bài viết. Khi bàn phím ảo xuất hiện làm thu hẹp không gian, \texttt{SingleChildScrollView} sẽ tự động kích hoạt thanh cuộn để người dùng kéo xem nội dung bị che khuất.
  \end{itemize}

  \vspace{0.25cm}
  \textbf{2. Phân chia không gian linh hoạt (\texttt{Expanded} hoặc \texttt{Flexible}):}
  \begin{itemize}
    \item \textbf{Nguyên tắc:} Thay vì để kích thước các widget tự do phình to gây đẩy các phần tử khác ra khỏi màn hình, ta ép chúng co giãn theo tỷ lệ không gian còn lại của trục chính (\texttt{Main Axis}) trong \texttt{Row} hoặc \texttt{Column}.
    \item \textbf{Áp dụng:}
      \begin{itemize}
        \item \texttt{Expanded} (Tight Fit): Ép widget con phải lấp đầy hoàn toàn phần không gian trống còn lại.
        \item \texttt{Flexible} (Loose Fit): Cho phép widget con co giãn linh hoạt nhưng không vượt quá kích thước tự nhiên tối đa của nó.
      \end{itemize}
  \end{itemize}

  \vspace{0.25cm}
  \textbf{3. Tự động bẻ xuống dòng (\texttt{Wrap}):}
  \begin{itemize}
    \item \textbf{Nguyên tắc:} Mặc định \texttt{Row} chỉ sắp xếp theo một hàng ngang duy nhất và sẽ gây Overflow nếu hết chỗ ngang. Widget \texttt{Wrap} giải quyết việc này bằng cách tự động ngắt dòng, đưa các phần tử không đủ chỗ xuống dòng tiếp theo (giống như cách chữ ngắt dòng trong văn bản Word).
    \item \textbf{Áp dụng:} Rất phù hợp cho danh sách các thẻ tag, bộ lọc tìm kiếm (\texttt{Chip}), hoặc danh sách từ khóa có độ dài không cố định.
  \end{itemize}
\end{frame}

% ==============================================================================
% TRANG 10: CƠ SỞ LÝ THUYẾT - 5 CÁCH XỬ LÝ OVERFLOW (PHẦN 2)
% ==============================================================================
\begin{frame}
  \SlideHeader{CƠ SỞ LÝ THUYẾT}{Chi Tiết Các Cách Xử Lý Overflow (Phần 2)}

  \begin{columns}[T]
    \begin{column}{0.52\textwidth}
      \textbf{4. Giới hạn hiển thị chữ dài (\texttt{TextOverflow.ellipsis} \& \texttt{maxLines}):}
      \begin{itemize}
        \item \textbf{Nguyên tắc:} Khi văn bản có độ dài quá lớn so với chiều rộng khung chứa, thay vì để tràn hoặc vỡ giao diện, ta giới hạn số dòng hiển thị và tự động thêm dấu ba chấm (\dots) vào cuối.
        \item \textbf{Áp dụng:} Dùng cho tiêu đề bài viết, tên sản phẩm, hoặc tin nhắn dài trong danh sách chat.
      \end{itemize}

      \vspace{0.25cm}
      \textbf{5. Co giãn tỷ lệ màn hình (\texttt{FittedBox} hoặc \texttt{MediaQuery}):}
      \begin{itemize}
        \item \textbf{Nguyên tắc:} Tự động thu nhỏ hoặc phóng to toàn bộ nội dung con cho vừa vặn với kích thước khung chứa của cha.
        \item \textbf{Áp dụng:} Thường dùng cho các icon, logo, hoặc tiêu đề cần vừa khít trên mọi kích cỡ màn hình.
      \end{itemize}
    \end{column}

    \begin{column}{0.44\textwidth}
      \centering
      \textbf{Cây quyết định lựa chọn giải pháp:}\\[0.2cm]
      \begin{tikzpicture}[font=\tiny, node distance=0.55cm]
        \node[draw=NavyTitle, fill=CardBg, rounded corners=4pt, line width=0.9pt] (root) {\textbf{Phát hiện nguy cơ tràn?}};
        
        \node[draw=AccentOrange, fill=AccentOrange!10, rounded corners=3pt, below left=0.45cm and -0.3cm of root] (horiz) {\textbf{Tràn Ngang (Row)}};
        \node[draw=SuccessGreen, fill=SuccessGreen!10, rounded corners=3pt, below right=0.45cm and -0.3cm of root] (vert) {\textbf{Tràn Dọc (Column)}};

        \draw[->, >=stealth, line width=0.8pt] (root) -- (horiz);
        \draw[->, >=stealth, line width=0.8pt] (root) -- (vert);

        \node[draw=gray!80, fill=white, rounded corners=3pt, below=0.35cm of horiz, text width=2.5cm, align=left] {
          $\bullet$ Văn bản dài: \textbf{Ellipsis}\\
          $\bullet$ Chia tỷ lệ: \textbf{Expanded}\\
          $\bullet$ Nhiều Tag: \textbf{Wrap}
        };

        \node[draw=gray!80, fill=white, rounded corners=3pt, below=0.35cm of vert, text width=2.5cm, align=left] {
          $\bullet$ Bàn phím: \textbf{ScrollView}\\
          $\bullet$ Đẩy chân: \textbf{Spacer}\\
          $\bullet$ Phóng/thu: \textbf{FittedBox}
        };
      \end{tikzpicture}
    \end{column}
  \end{columns}
\end{frame}

% ==============================================================================
% TRANG 11: YÊU CẦU THỰC HÀNH - THIẾT KẾ CẤU TRÚC 7 THÀNH PHẦN
% ==============================================================================
\begin{frame}[fragile]
  \SlideHeader{YÊU CẦU THỰC HÀNH}{Thiết Kế Cấu Trúc Bố Cục Dự Án Demo 7 Thành Phần}

  \begin{columns}[T]
    \begin{column}{0.52\textwidth}
      \textbf{Cây thư mục mã nguồn (\texttt{chu\_de\_tuan12/lib/})}
\begin{lstlisting}[numbers=none, basicstyle=\ttfamily\tiny]
lib/
|-- screens/login/
|   |-- login_screen.dart       # Dieu phoi bo cuc goc
|   \-- widgets/
|       |-- login_header.dart   # 1. Logo & Slogan
|       |-- account_card.dart   # 2. Avatar & Stack Badge
|       |-- password_field.dart # 3. Input mat khau
|       |-- login_button.dart   # 4. Nut Expanded chinh
|       |-- face_id_button.dart # 5. Nut Fixed 52x52
|       |-- login_options_row.dart # 6. Row spaceBetween
|       \-- login_footer.dart   # 7. Spacer day xuong day
\end{lstlisting}
    \end{column}

    \begin{column}{0.46\textwidth}
      \textbf{Phân bổ trách nhiệm không gian:}
      \begin{itemize}
        \item \textbf{1. Header:} Column căn giữa trục phụ.
        \item \textbf{2. AccountCard:} Phối hợp Row, Column, Stack và Expanded chống tràn tên dài.
        \item \textbf{3. PasswordField:} Ô nhập co dãn ngang.
        \item \textbf{4 \& 5. Button Row:} Expanded (flex: 4) kết hợp Nút cố định $52 \times 52\text{px}$.
        \item \textbf{6. Options Row:} Phân bổ 2 đầu bằng \texttt{spaceBetween}.
        \item \textbf{7. Footer:} Ghim chân trang qua \texttt{Spacer()}.
      \end{itemize}
    \end{column}
  \end{columns}
\end{frame}

% ==============================================================================
% TRANG 12: MINH HỌA MÃ NGUỒN TÁC GIẢ - CĂN CHỈNH ROW & COLUMN
% ==============================================================================
\begin{frame}[fragile]
  \SlideHeader{MINH HỌA MÃ NGUỒN TÁC GIẢ}{Căn Chỉnh Trục: MainAxisAlignment \& CrossAxisAlignment}

  \begin{columns}[T]
    \begin{column}{0.48\textwidth}
      \textbf{\texttt{widgets/login\_options\_row.dart}}
\begin{lstlisting}
class LoginOptionsRow extends StatelessWidget {
  @override
  Widget build(BuildContext context) {
    return Row(
      // 1. Day 2 tuy chon ve 2 mep trai va phai
      mainAxisAlignment: 
          MainAxisAlignment.spaceBetween,
      // 2. Can giua truc doc de chu thang hang
      crossAxisAlignment: 
          CrossAxisAlignment.center,
      children: [
        Flexible(
          child: TextButton(
            onPressed: onForgotPassword,
            child: const Text('Quen mat khau?'),
          ),
        ),
        const SizedBox(width: 12),
        Flexible(
          child: TextButton.icon(
            onPressed: onActivateOtp,
            icon: const Icon(Icons.shield_outlined),
            label: const Text('Kich hoat Smart OTP'),
          ),
        ),
      ],
    );
  }
}
\end{lstlisting}
    \end{column}

    \begin{column}{0.48\textwidth}
      \textbf{\texttt{screens/login/login\_screen.dart}}
\begin{lstlisting}
// Column chinh cua man hinh dang nhap:
Column(
  // 1. Main Axis (Doc): Bat dau tu dinh
  mainAxisAlignment: MainAxisAlignment.start,
  // 2. Cross Axis (Ngang): Keo gian deu cac 
  // widget con de khop voi chieu ngang Form
  crossAxisAlignment: CrossAxisAlignment.stretch,
  children: [
    const LoginHeader(),
    const SizedBox(height: 24),
    AccountCard(...),
    const SizedBox(height: 18),
    PasswordField(...),
    const SizedBox(height: 20),
    // Hang nut Dang nhap & Face ID
    Row(...),
    const SizedBox(height: 14),
    const LoginOptionsRow(),
    const Spacer(), // Day footer sat day
    const LoginFooter(),
  ],
)
\end{lstlisting}
    \end{column}
  \end{columns}
\end{frame}

% ==============================================================================
% TRANG 13: MINH HỌA MÃ NGUỒN TÁC GIẢ - EXPANDED & NÚT CỐ ĐỊNH
% ==============================================================================
\begin{frame}[fragile]
  \SlideHeader{MINH HỌA MÃ NGUỒN TÁC GIẢ}{Phân Bổ Không Gian Hàng Nút Bằng Expanded \& Fixed}

  \textbf{\texttt{screens/login/login\_screen.dart (Hàng nút Đăng nhập \& Face ID)}}
\begin{lstlisting}
// Hang nut hanh dong: Ket hop giua nut co gian linh hoat va nut co dinh
Row(
  crossAxisAlignment: CrossAxisAlignment.center,
  children: [
    // 1. Nut Dang Nhap boc Expanded: Tu dong chiem tron toan bo khong gian con lai
    Expanded(
      child: LoginButton(
        isLoading: _isLoading,
        onPressed: _handleLogin,
      ),
    ),
    const SizedBox(width: 12), // Khoang dem co dinh giua 2 nut
    // 2. Nut Face ID giu nguyen kich thuoc hinh vuong co dinh 52x52 pixel
    const FaceIdButton(), 
  ],
)
\end{lstlisting}

  \begin{tcolorbox}[colback=CardBg, colframe=NavyHeader!30, arc=4pt, boxrule=0.8pt, top=5pt, bottom=5pt]
    {\scriptsize \textbf{Phân tích kỹ thuật của tác giả trong bài:} Nút Đăng nhập được bọc trong \texttt{Expanded} sẽ tự động giãn rộng để vừa khít từ màn hình nhỏ nhất (320px) đến máy tính bảng, trong khi nút Face ID bên cạnh luôn giữ vững tỷ lệ $52 \times 52\text{px}$ vuông vức hoàn hảo với hiệu ứng khung quét \texttt{Stack \& Positioned} mà không bị bóp méo.}
  \end{tcolorbox}
\end{frame}

% ==============================================================================
% TRANG 14: NHẤN MẠNH CÁCH XỬ LÝ CỦA TÔI - TRÀN NGANG (ELLIPSIS)
% ==============================================================================
\begin{frame}[fragile]
  \SlideHeader{NHẤN MẠNH GIẢI PHÁP CỦA TÔI}{Trị Dứt Điểm Tràn Ngang Văn Bản: Expanded + Ellipsis}

  \begin{columns}[T]
    \begin{column}{0.48\textwidth}
      \textbf{\textcolor{ErrorRed}{Lỗi phổ biến (RenderFlex Overflowed):}}
\begin{lstlisting}[rulecolor=\color{ErrorRed}]
Row(
  children: [
    AvatarWidget(),
    // KHONG boc Expanded!
    Column(
      children: [
        Text(
          'NGUYEN HOANG MINH KHANG TRONG DAT',
          // Text vuot mep phai man hinh!
        ),
      ],
    ),
  ],
)
// Hau qua: A RenderFlex overflowed by 142px
\end{lstlisting}
    \end{column}

    \begin{column}{0.48\textwidth}
      \textbf{\textcolor{SuccessGreen}{Đoạn code xử lý của tôi trong AccountCard:}}
\begin{lstlisting}[rulecolor=\color{SuccessGreen}]
// File: widgets/account_card.dart
Expanded(
  child: Column(
    crossAxisAlignment: CrossAxisAlignment.start,
    children: [
      Text(
        accountName,
        maxLines: 1, // Gioi han 1 dong
        // Tu dong rut gon dau ba cham an toan
        overflow: TextOverflow.ellipsis,
        style: const TextStyle(
          fontSize: 15, fontWeight: FontWeight.bold),
      ),
      const SizedBox(height: 3),
      Text(accountNumber, style: ...),
    ],
  ),
)
\end{lstlisting}
    \end{column}
  \end{columns}

  \vspace{0.1cm}
  \begin{tcolorbox}[colback=SuccessGreen!5, colframe=SuccessGreen, arc=4pt, boxrule=0.8pt, top=3pt, bottom=3pt]
    {\scriptsize \textbf{Minh chứng thực nghiệm trong bài:}\\
    Tôi đã tích hợp nút \textbf{"Đổi tài khoản"} ngay trên thẻ \texttt{AccountCard}. Khi bấm, ứng dụng hoán đổi giữa tên ngắn chuẩn sang tên siêu dài (\texttt{AppStrings.longAccountName}: 65 ký tự). \texttt{Expanded} ép \texttt{Text} nhận đúng độ rộng khả dụng và lập tức hiển thị \texttt{NGUYỄN HOÀNG MINH KHANG...} mà hoàn toàn không bị lỗi vạch sọc vàng-đen!}
  \end{tcolorbox}
\end{frame}

% ==============================================================================
% TRANG 15: NHẤN MẠNH CÁCH XỬ LÝ CỦA TÔI - BỘ 4 TRỊ TRÀN DỌC
% ==============================================================================
\begin{frame}[fragile]
  \SlideHeader{NHẤN MẠNH GIẢI PHÁP CỦA TÔI}{Bộ 4 Kinh Điển Trị Tràn Dọc Bàn Phím Ảo}

  \textbf{\texttt{screens/login/login\_screen.dart (Bộ điều phối không gian thích ứng triệt để)}}
\begin{lstlisting}
LayoutBuilder(builder: (context, constraints) {
  return Center(
    child: ConstrainedBox(
      constraints: BoxConstraints(maxWidth: maxFormWidth), // 1. Khong che chieu rong
      child: SingleChildScrollView( // 2. Tu dong kich hoat cuon khi chieu cao bi ep chat
        physics: const BouncingScrollPhysics(),
        child: ConstrainedBox(
          // 3. Ep chieu cao toi thieu bang chieu cao man hinh de giu chan Footer
          constraints: BoxConstraints(minHeight: constraints.maxHeight - 32.0),
          child: IntrinsicHeight( // 4. Tinh toan chieu cao thuc te de Spacer() hoat dong
            child: Column(
              crossAxisAlignment: CrossAxisAlignment.stretch,
              children: [
                const LoginHeader(),
                AccountCard(...),
                PasswordField(...), // Khi bam vao -> Ban phim ao day len!
                Row(...),
                const LoginOptionsRow(),
                const Spacer(), // Khi man hinh du cao: Day Footer sat xuong day!
                const LoginFooter(),
              ],
            ),
          ),
        ),
      ),
    ),
  );
});
\end{lstlisting}

  \begin{tcolorbox}[colback=CardBg, colframe=NavyHeader!30, arc=4pt, boxrule=0.8pt, top=3pt, bottom=3pt]
    {\scriptsize \textbf{Phân tích cơ chế kép độc đáo:}\\
    $\bullet$ \textbf{Khi màn hình bình thường:} \texttt{Spacer()} bung ra hết cỡ đẩy Footer xuống sát mép đáy trang trọng.\\
    $\bullet$ \textbf{Khi mở bàn phím ảo hoặc xoay ngang máy:} Chiều cao khả dụng bị co hẹp $\rightarrow$ \texttt{SingleChildScrollView} tự động nhận diện và cho phép cuộn mượt mà, triệt tiêu 100\% lỗi tràn dọc (Bottom Overflow)!}
  \end{tcolorbox}
\end{frame}

% ==============================================================================
% TRANG 16: NHẤN MẠNH CÁCH XỬ LÝ CỦA TÔI - FLEXIBLE & RESPONSIVE
% ==============================================================================
\begin{frame}[fragile]
  \SlideHeader{NHẤN MẠNH GIẢI PHÁP CỦA TÔI}{Chống Tràn Với Flexible(fit: loose) \& Responsive}

  \begin{columns}[T]
    \begin{column}{0.48\textwidth}
      \textbf{1. Áp dụng \texttt{Flexible(fit: loose)}}\\
      \textbf{\texttt{widgets/login\_options\_row.dart}}
\begin{lstlisting}
// Thay vi de TextButton tu do tran hang:
Flexible(
  fit: FlexFit.loose, // Co gian tu nhien
  child: TextButton(
    onPressed: onForgotPassword,
    child: const Text(
      'Quen mat khau?',
      style: TextStyle(fontSize: 13),
    ),
  ),
),
\end{lstlisting}
      {\scriptsize $\rightarrow$ Khi hệ điều hành phóng to font chữ (\textit{Text Scale Factor}), \texttt{Flexible} đảm bảo 2 nút tùy chọn tự co ép trong biên độ cho phép mà không làm tràn ngang.}
    \end{column}

    \begin{column}{0.48\textwidth}
      \textbf{2. Khống chế Breakpoint đa kích thước}\\
      \textbf{\texttt{screens/login/login\_screen.dart}}
\begin{lstlisting}
final bool isWideScreen = 
    constraints.maxWidth > 600;
final double horizontalPadding = 
    isWideScreen ? 32.0 : 20.0;
// Tablet / Desktop / Web: Max 440px can giua
final double maxFormWidth = 
    isWideScreen ? 440.0 : double.infinity;

ConstrainedBox(
  constraints: BoxConstraints(
    maxWidth: maxFormWidth),
  child: ...,
)
\end{lstlisting}
      {\scriptsize $\rightarrow$ Giữ cho form đăng nhập luôn cân đối, chuyên nghiệp trên máy tính bảng và màn hình PC rộng.}
    \end{column}
  \end{columns}
\end{frame}

% ==============================================================================
% TRANG 17: PHÂN TÍCH CHUYÊN SÂU - BẢNG SO SÁNH KỸ THUẬT
% ==============================================================================
\begin{frame}
  \SlideHeader{PHÂN TÍCH CHUYÊN SÂU}{Bảng So Sánh Toàn Diện Các Giải Pháp Xử Lý Overflow}

  \centering
  \renewcommand{\arraystretch}{1.2}
  {\scriptsize
  \begin{tabular}{p{2.6cm}p{1.9cm}p{2.3cm}p{2.5cm}p{3.3cm}}
    \toprule
    \rowcolor{TableHeadBg}
    \textbf{Giải pháp} & \textbf{Chiều Áp Dụng} & \textbf{Hành Vi Kích Thước} & \textbf{Khả Năng Xử Lý Tràn} & \textbf{Áp Dụng Trong Dự Án} \\
    \midrule
    \textbf{SingleChildScrollView} & Trục dọc / ngang & Mở rộng vô hạn trục cuộn & Tự kích hoạt thanh cuộn & Bọc toàn bộ Form chống tràn bàn phím \\
    \rowcolor{TableAltBg}
    \textbf{Expanded (Tight)} & Trục chính Row/Col & Ép lấp đầy (\textit{Tight fit}) & Không thể tự tràn & Nút LoginButton \& AccountCard \\
    \textbf{Flexible (Loose)} & Trục chính Row/Col & Tự do $\le$ max (\textit{Loose fit}) & Giữ size tự nhiên $\le$ trần & Hàng tùy chọn LoginOptionsRow \\
    \rowcolor{TableAltBg}
    \textbf{Wrap} & Đa dòng (2D) & Tự bẻ xuống dòng mới & Trị dứt điểm tràn ngang & Danh sách thẻ tag, bộ lọc chip \\
    \textbf{Ellipsis \& MaxLines} & Trục chữ (Text) & Giới hạn số dòng cố định & Cắt chuỗi và thêm dấu \dots & Tên tài khoản trong AccountCard \\
    \rowcolor{TableAltBg}
    \textbf{FittedBox / Scale} & 2 chiều (XY) & Thu nhỏ tỷ lệ cả khung & Thu nhỏ widget vừa khung & Icon, Logo và nhãn cố định \\
    \bottomrule
  \end{tabular}
  }
\end{frame}

% ==============================================================================
% TRANG 18: KIỂM CHỨNG & THỰC NGHIỆM - DEMO THỰC TẾ
% ==============================================================================
\begin{frame}
  \SlideHeader{KIỂM CHỨNG \& THỰC NGHIỆM}{Kiểm Thử Thực Tế Các Kịch Bản Thay Đổi Không Gian}

  \begin{columns}[T]
    \begin{column}{0.48\textwidth}
      \textbf{1. Kịch bản Nội dung thay đổi (Content Mutation):}
      \begin{itemize}
        \item Bấm nút icon \textit{Swap} trên thẻ \texttt{AccountCard}.
        \item Tên tài khoản chuyển sang chuỗi siêu dài gồm 65 ký tự tiếng Việt có dấu.
        \item \textbf{Kết quả:} \texttt{Expanded} lập tức giới hạn độ rộng, chuỗi tự thu gọn thành \texttt{NGUYỄN HOÀNG MINH KHANG TRỌNG ĐẠT...}, giao diện hoàn toàn nguyên vẹn 100\%.
      \end{itemize}

      \vspace{0.15cm}
      \textbf{2. Kịch bản Xoay ngang màn hình (Orientation Change):}
      \begin{itemize}
        \item Chiều cao giảm từ 800px xuống còn 400px.
        \item \textbf{Kết quả:} \texttt{SingleChildScrollView} tự kích hoạt, cho phép vuốt cuộn êm ái.
      \end{itemize}
    \end{column}

    \begin{column}{0.48\textwidth}
      \textbf{3. Kịch bản Bàn phím ảo đẩy lên (Virtual Keyboard):}
      \begin{itemize}
        \item Chạm vào ô nhập mật khẩu \texttt{PasswordField}.
        \item Bàn phím chiếm 280px chiều cao màn hình.
        \item \textbf{Kết quả:} \texttt{Spacer()} tự co lại, toàn bộ các nút bấm và footer cuộn lên tự nhiên, không xuất hiện vạch sọc vàng-đen.
      \end{itemize}

      \vspace{0.15cm}
      \textbf{4. Kịch bản Kích thước lớn (Tablet / Web Desktop):}
      \begin{itemize}
        \item Chiều rộng màn hình $> 600\text{px}$.
        \item \textbf{Kết quả:} \texttt{LayoutBuilder} tự khống chế \texttt{maxWidth: 440px} và căn giữa hoàn mỹ.
      \end{itemize}
    \end{column}
  \end{columns}
\end{frame}

% ==============================================================================
% TRANG 19: TỔNG KẾT ĐỀ TÀI - 4 TRỤ CỘT CỐT LÕI ĐÃ ĐẠT ĐƯỢC
% ==============================================================================
\begin{frame}
  \SlideHeader{TỔNG KẾT ĐỀ TÀI}{4 Trụ Cột Bố Cục Cốt Lõi Đã Đạt Được}

  \textbf{1. Nắm vững bản chất nguyên lý Box Constraints \& Trục toạ độ}\\
  {\small Thấu suốt quy tắc \textit{"Constraints go down, Sizes go up, Parent sets position"}, làm chủ hướng trục Main/Cross Axis để điều phối không gian chính xác.}\\[0.22cm]

  \textbf{2. Nắm trọn vẹn 5 phương pháp xử lý Overflow chuẩn công nghiệp}\\
  {\small Hiểu sâu nguyên lý và kịch bản áp dụng của: Cuộn trang (\texttt{SingleChildScrollView}), Phân chia linh hoạt (\texttt{Expanded/Flexible}), Tự bẻ dòng (\texttt{Wrap}), Giới hạn chữ (\texttt{Ellipsis}), và Co giãn tỷ lệ (\texttt{FittedBox}).}\\[0.22cm]

  \textbf{3. Đóng gói giải pháp thực tế xuất sắc trong mã nguồn dự án}\\
  {\small Trị dứt điểm tràn ngang bằng combo \texttt{Expanded + TextOverflow.ellipsis} và trị dứt điểm tràn dọc bằng bộ 4 \texttt{LayoutBuilder + ScrollView + ConstrainedBox + IntrinsicHeight + Spacer}}.\\[0.22cm]

  \textbf{4. Xây dựng ứng dụng mẫu Responsive \& Adaptive hoàn chỉnh}\\
  {\small Bố cục 7 widget thành phần tự động co giãn và thích ứng hoàn hảo trên mọi kích thước màn hình cũng như khi nội dung thay đổi.}
\end{frame}

% ==============================================================================
% TRANG 20: KẾT THÚC BÁO CÁO - CẢM ƠN & Q&A
% ==============================================================================
\begin{frame}
  \vspace{0.7cm}
  \centering
  {\textcolor{AccentOrange}{\textbf{\footnotesize KẾT THÚC BÁO CÁO}}}\\[0.35cm]
  {\Huge \textbf{\textcolor{NavyTitle}{XIN CHÂN THÀNH CẢM ƠN!}}}\\[0.25cm]
  {\Large \textbf{\textcolor{TextDark}{Thầy Cô \& Các Bạn Đã Chú Ý Lắng Nghe}}}\\[0.5cm]

  \begin{tcolorbox}[colback=CardBg, colframe=NavyHeader!40, width=8.5cm, arc=8pt, boxrule=1.1pt, halign=center, top=8pt, bottom=8pt]
    {\Huge \textbf{\textcolor{NavyTitle}{Q \& A}}}\\[0.15cm]
    {\small \textbf{Phiên thảo luận và giải đáp thắc mắc chuyên đề kỹ thuật}}\\[0.1cm]
    {\footnotesize Chủ đề 02: Tổ chức bố cục linh hoạt và xử lý không gian trong Flutter}
  \end{tcolorbox}
\end{frame}

\end{document}
